%%%%%%%%%%%%%%%%%%%%%%%%%%%%%%%%%%%%%%%%%%%%%%%%%%%%%%%%%
% --------------------- INTRODUCTION --------------------
%%%%%%%%%%%%%%%%%%%%%%%%%%%%%%%%%%%%%%%%%%%%%%%%%%%%%%%%%
\section{Introduction}\label{sec:introduction}

%\manolis{I commented some parts}, you can find them by searching \verb|\begin{comment}|.

%\manolis{General remark: when citing multiple papers, respect the order, e.g., [2,4,35] instead of [2,35,4].}

%\manolis{I throw stuff in the appendix if I believe that they are fine as is (with the exception of the DPs).
%Of course, feel free to read and comment on them.}



    The fair division of indivisible goods is a fundamental problem in computer science, with numerous applications, such as distribution of inheritance, scheduling jobs, and allocating computational resources to machine learning models. The concept of fair division of \emph{divisible} resources was introduced by Steihaus~\cite{h__steihaus_1948}; in recent years, interest has shifted towards \emph{indivisible} goods, due both to the aforementioned applications and to several intriguing open questions surrounding the study of \emph{envy}. Since envy-free (EF) allocations are not guaranteed to exist, research has focused on \emph{envy minimization} (e.g.,~\cite{ISAAC_envy_ratio,LiptonEtAl,Nguyen_Rothe}),
    and \emph{relaxations} of EF, such as EF1 \cite{budish} and EFX \cite{EFXCara,Laurent_EFX}. The existence of EFX allocations remains a major open question in fair division (cf.~\cite{FD_survey}),
    with some results for two~\cite{PlautRough} and three agents~\cite{akrami2022efxallocationssimplificationsimprovements,CGM24}, and further work limiting agents' valuations~\cite{TwoValuedInstanses,babaioff2020fairtruthfulmechanismsdichotomous,DBLP:conf/aaai/HosseiniSVX21,efxexistsfor3typesofadditiveagents}.

    A restricted version of the problem was introduced by Christodoulou, Fiat, Koutsoupias and Sgouritsa~\cite{EFXgraphs}, with a graph model where the vertices represent agents and the edges represent goods; it was proven in the same paper that EFX allocations always exist in this restricted version for simple graphs. This model has been extensively studied since, with the existence of EFX allocations also being studied in multigraphs \cite{AfshinmehretTriangleFree,AfshinmehrDKMR25,bhaskar_et_al:LIPIcs.FSTTCS.2025.15_efxallocationsmultigraphclasses,kaviani2024envyfreeallocationindivisiblegoods,OntheExistenceofEFXAllocationsinMultigraphs}. An interesting concept in this model is that of \emph{EFX orientations}, i.e., EFX allocations where edges have to be allocated to one of their endpoints, naturally modeling practical applications where resources are contested only by neighbors. EFX orientations have also received significant attention in recent years, with their parameterized complexity in particular being studied in \cite{AfshinmehrDKMR25,sagt/BlazejGRS25,ijcai/DeligkasEGK25}.

    In spite of this, EF orientations remain largely unexplored.
    In this work we study them primarily under the lens of parameterized complexity, and present a plethora of tractable special cases, parameterized algorithms, and hardness results, including a complete characterization for the case of binary valuations (\cref{theorem:connected_component_props}). Many of our results transfer to EFX orientations, thus complementing and improving upon results from previous work on EFX orientations~\cite{AfshinmehrDKMR25,sagt/BlazejGRS25,ijcai/DeligkasEGK25}; see \cref{sec:contributions} for details.

    Another interesting relaxation of EF(X) that we consider in this work concerns allocations in which not every item is allocated to an agent, i.e., EF(X) with \emph{charity} \cite{CGH19}.
    To the best of our knowledge, we are the first to consider charity in the graph setting of \cite{EFXgraphs}, establishing algorithms for finding the minimum number of edges that have to remain unoriented (i.e., donated to charity) in order for EF(X) orientations to exist.

    %A different restricted instance was introduced in \cite{EFXgraphs} to tackle the problem. Agents are represented as vertices of a graph and goods are represented as edges of the same graph. Only endpoints of an edge have any value for that edge. It was shown that an EFX allocation always exists in graphs. This special instance is equivalent to the general problem when considering complete multi-hypergraphs. The existence of EFX allocations in multigraphs has been explored recently \cite{AfshinmehrDKMR25,bhaskar_et_al:LIPIcs.FSTTCS.2025.15_efxallocationsmultigraphclasses,kaviani2024envyfreeallocationindivisiblegoods,OntheExistenceofEFXAllocationsinMultigraphs} \MinasComment{Kaviani et Al paper is in WINE2024 BUT there was an issue with the proceedings, so it exists only in Arxiv}. In \cite{EFXgraphs}, it was also shown that EFX orientation do not need to always exists by providing a counterexample of a clique of size 4. An EFX orientation is an EFX allocation in which each good is given to one of the endpoints. In graphs and multigraphs there may be EFX allocations that are not EFX orientations. These kinds of allocations can be wasteful, in the sense that items may be allocated to agents that do not value them. An interesting question is: Can we find EFX orientations in multigraphs efficiently? It is important to mention that EF1 orientations always exist in graphs and hypergraphs \cite{EFXgraphs,ijcai/DeligkasEGK25}.
    %In \cite{ifaamas/ZengM25} they characterized when given a graph if there are certain valuations over the edges such that the graph has an EFX orientation. In \cite{sagt/BlazejGRS25,ijcai/DeligkasEGK25} they explored how to compute EFX orientations used different parameters based on the multigraph. We generalize a result in \cite{sagt/BlazejGRS25} for multigraphs regarding the parameterized.


%Fair division is a problem relevant to this day, allocation of resources between different parties appears constantly in life, for instance rent division between roommates, distribution of inheritance, and for more technological examples, consider scheduling jobs between servers and computational resources between machine learning models.
    %Fair division was first introduced in \cite{h__steihaus_1948}.


    % Later on the fairness criterion of envy-freeness was introduced and studied, where agents do not prefer what others agents have, \cite{foley1966resource,gamow1958puzzle,VARIAN197463}. This type of fair allocations always exist for divisible items \cite{Aziz2016,Stromquist1980HowTC,Woo80}. Regarding indivisible items, this fairness criterion is too strict, simply consider two agents who both value an indivisible good, whoever gets the item is going to be envied.


 %\begin{description}

  %\item[Question:] Can we extend the FPT algorithm of \cite{ijcai/DeligkasEGK25} to other related problems? For example: \emph{What is the minimum amount of edges that have to be removed from a given graph, in order for an EFX orientation to exist?} (The concept of "removing goods" is common in fair division, where it is known as \emph{charity}. It is used to eliminate or reduce envy.)

  %\item[Summary:] EFX orientations are a new variation of EFX allocations, introduced in \cite{sigecom/CFKS23}. Imagine a graph where the edges are goods and the nodes are agents. Each edge $(u,v)$ (good) has to be oriented towards $u$ or $v$, i.e. allocated to one of two agents. Nodes $u$ and $v$ (agents) might have different utilities for the same edge, but the edge $(u,v)$ has zero utility for everyone else. If the edges can be oriented in a way that no agent envies any other \emph{up to any good}, then this is an \emph{EFX orientation}.

  %From a practical perspective, this problem makes sense in contexts where each resource is wanted only by two specific agents, for example borders between countries.

  %There are many recent works on EFX orientations, e.g. \cite{ifaamas/ZengM25} which relates the existence of EFX orientations to chromatic numbers and \cite{ijcai/ZWL024} which studies the case where the edges may be \emph{chores} instead of \emph{goods}. The concept of EFX orientations has sparked great interest, due to some existence results which we lack for regular EFX allocations.

  %\item[People interested:] Sotiris, Christos, Manolis, Michael, Laurent, Aris

  %\item[Comments:]
  %Sotiris:

  %Manolis: A case that remains open is the parameterization by the distance to clique.

  %Michael: It seems the hardness results for vertex cover or bounded-order multipgraphs are only \emph{weak} NP-hardness. Are pseudopolynomial time algorithms possible? In particular, is the problem in P for constant treewidth and unary weights?


  %Laurent: what about fairness concepts other than EF? Proportionality for example (every agent wants to have utility at least $1/n$ where $n$ is the number of agents).
%\end{description}

\subsection{Related Work}

\paragraph{Allocations.} Regarding EF allocations, NP-completeness and inapproximability results are shown by Lipton, Markakis, Mossel and Saberi~\cite{LiptonEtAl}.
Recent results for the existence of EFX allocations include the following: Plaut and Roughgarden~\cite{PlautRough} show existence for two agents with general monotone valuations; Chaudhury, Garg and Mehlhorn~\cite{CGM24} for three agents with additive valuations; Akrami et al.~\cite{akrami2022efxallocationssimplificationsimprovements} improve this to a slightly more general class; Prakash, Ghosal, Nimbhorkar and Varma~\cite{efxexistsfor3typesofadditiveagents} show existence when agents have one of three additive valuations.



\paragraph{Orientations.}
The existence of EFX orientations is known to be NP-complete even for symmetric valuations by Christodoulou et al.~\cite{EFXgraphs}.
Deligkas, Eiben, Goldsmith and Korchemna~\cite{ijcai/DeligkasEGK25} prove that the problem remains NP-complete even for graphs with vertex cover size $8$ or multigraphs with $10$ vertices, while also identifying one of the few known fixed-parameter tractable cases through the parameterization by \emph{slim tree-cut width} (defined in \cite{algorithmica/GanianK24}).
%They also mention that the problem is NP-complete for graphs of bounded maximum degree by the proof of~\cite{EFXgraphs}.
Blazej, Gupta, Ramanujan and Strulo~\cite{sagt/BlazejGRS25} study the complexity of the problem in relation to how close the graph is to being bipartite,
proving that it remains NP-complete for (a subclass of) binary instances,
while also showing a linear-time algorithm for symmetric binary instances of constant treewidth.
Afshinmehr, Danaei, Kazemi, Mehlhorn and Rathi~\cite{AfshinmehrDKMR25} and Hsu~\cite{ecai/Hsu25} further study EFX orientations in (bipartite) multigraphs. %, with the latter in particular showing that the problem is NP-complete even with $8$ vertices.
Zeng and Mehta~\cite{ifaamas/ZengM25} relate the existence of EFX orientations to the graph's chromatic number.
EFX orientations have also been studied in the setting of chores, i.e., negatively valued items, and mixed manna, i.e., goods and chores \cite{hsu2025polynomialtimealgorithmefxorientationsChores,ijcai/ZWL024}.
Despite all the aforementioned work for EFX orientations, to the best of our knowledge,
the only work relevant to EF orientations appears in \cite{aldt/MisraS24},
where the authors study EF \emph{allocations} in simple graphs,
with a polynomial-time algorithm for binary valuations also transferring to \emph{orientations}.



\paragraph{Charity.}
The concept of charity was first introduced by Caragiannis, Gravin, and Huang~\cite{CGH19}, in the setting of EFX allocations.
Chaudhury, Kavitha, Mehlhorn and Sgouritsa~\cite{CKMS21} proved that the number of unallocated items that guarantees existence of EFX for general monotone valuations is upper bounded by $n-1$, where $n$ is the number of agents. This bound was later reduced to $n-2$ for additive valuations \cite{AlmostFullEFXforFourAgents}, and extended to general monotone valuations~\cite{mahara2021extensionadditivevaluationsgeneral}. For 4 agents, this number is 1 \cite{AlmostFullEFXforFourAgents}.
To the best of our knowledge, charity has never been studied in the context of fair orientations prior to this work, although one could obtain a simple charity bound of $|E|/2$ for EFX orientations in simple graphs by removing edges to obtain a bipartite graph~\cite{erd/Erdos65},
which is known to always admit an EFX orientation \cite{ifaamas/ZengM25}.

%\EdouardComment{I'm not sure if we already talked about this, but I am pretty sure the FPT algorithm parameterized by treewidth can be adapted to hypergraphs (with a worse complexity though). Maybe we can include it if we have time.}



\subsection{Our Contributions}\label{sec:contributions}



We start by considering in \cref{subsec:EFo:linear} the {\EFo} problem%
\footnote{For formal definitions of all the problems in question we refer to \cref{sec:preliminaries}.} in multigraphs with binary valuation functions.
We provide a complete characterization of EF orientations in this case (\cref{theorem:connected_component_props}), roughly stating that an EF orientation exists if and only if every connected component of the graph satisfies at least one among four graph properties. Our characterization implies a linear-time algorithm for {\EFo} in the case of binary valuations; in fact, the algorithm can be extended to the more general {\EFmc} problem (i.e., finding the minimum amount of edges that have to be removed for an EF orientation to exist).
%Towards obtaining this algorithm we prove a structural theorem (\cref{theorem:connected_component_props}) that roughly says that an EF orientation exists if and only if each of the graph's connected components has at least one among four properties.
%Then, we check each condition in linear time, thus resulting in the claimed running time.
This result contrasts the NP-hardness of {\EFXo} in the analogous setting,
which persists even for symmetric instances on simple graphs that are two edge-deletions away from being bipartite~\cite{sagt/BlazejGRS25}. It also generalizes to multigraphs the algorithm of Misra and Sethia~\cite{aldt/MisraS24} for \EFo\ in simple graphs.

In \cref{subsec:EF:hardness} we present various hardness results for {\EFo}.
First we show that the problem remains NP-hard even on $3$-regular simple graphs with weights in $\{1,2\}$ and symmetric valuations, through a reduction from an appropriate variant of SAT.
Given this hardness for very restrictive valuations,
we next ask whether restricting the \emph{structure} of the input graph allows for more efficient algorithms.
Via a simple reduction from \textsc{Partition} we show that, for arbitrary weights, %that is not the case and
{\EFo} remains NP-hard even for simple graphs of vertex cover number~$2$ and multigraphs with~$2$ vertices.
Therefore, we move on to consider instances with polynomially-bounded weights, in which case
we prove W[1]-, XNLP-, and XALP-hardness for the parameterization by vertex cover, pathwidth, and treewidth respectively (\cref{thm:EF:w1h}) through a reduction from the {\TOO} problem~\cite{wg/BodlaenderCW22,dam/BodlaenderS26}. We thus argue that the problem remains highly intractable even in structured inputs.
As a byproduct of our reduction in \cref{thm:EF:w1h} and in conjunction with a recent result on the fine-grained complexity of \UBP~\cite{bringmann2026tightsethbasedlowerbounds}, we obtain a $n^{o(\vc)}$ ETH-based lower bound,
where $\vc$ denotes the vertex cover number of the input graph, that holds even for simple graphs.

Although the hardness results of \cref{subsec:EF:hardness} seem very discouraging,
we next identify restricted cases on which we are able to obtain fast algorithms.
First, we consider instances on simple graphs with at most $k$ \emph{heavy} edges, that is, edges that have valuation at least $2$ for an endpoint;
this naturally generalizes the binary-valuation regime considered in \cref{subsec:EFo:linear}, albeit restricted to simple graphs.
Here we develop a $2^k n^{\bO(1)}$ algorithm, which we complement with a tight SETH-based lower bound (\cref{thm:seth-ef,thm:EFo:Heavy}).
We view this as a first step towards generalizing the algorithm of \cref{thm:EFo:linear_algo} to an FPT algorithm parameterized by the
number of heavy edges. %, which constitutes a natural metric of ``distance to binary valuation''
Second, using dynamic programming (DP), we show that {\EFo} can be solved in time $(n+m)^{\bO(\tw)}$ for polynomially-bounded weights,
where $\tw$ denotes the treewidth of the graph and $m$ the number of edges.
Our algorithm works for the more general {\EFmc} problem in multigraphs with non-symmetric valuations,
and matches the lower bound of \cref{thm:EF:w1h}.

{\EFo} is arguably a very natural, understudied problem and our results, apart from being interesting in their own right,
additionally pave the way to advance the state-of-the-art for the closely related, and better understood, {\EFXo} problem.
In particular, in \cref{subsec:EF->EFX} we present a reduction from {\EFo} to {\EFXo} that allows us to transfer most hardness results of \cref{subsec:EF:hardness}
to the latter. Consequently, we obtain NP-hardness for simple graphs
of vertex cover number~$4$ and multigraphs of~$4$ vertices, improving upon the hardness results of Deligkas et al.~\cite{ijcai/DeligkasEGK25} for constants~$8$ and~$10$ respectively, as well as a hardness result by Afshinmehr et al.~\cite{AfshinmehrDKMR25} for multigraphs with~$8$ vertices.
More importantly, we obtain a $n^{o(\vc)}$ ETH-based lower bound for simple symmetric instances with polynomially-bounded weights.
As a matter of fact, in \cref{subsec:EFXo:tw} we yet again employ DP and obtain an $(n+m)^{\bO(\tw)}$ algorithm,
this time for the {\EFXmc} problem; once again, our algorithm works for multigraphs and non-symmetric valuations.
Since this matches the aforementioned lower bound, we settle the open question of Deligkas et al.~\cite{ijcai/DeligkasEGK25} on the complexity of {\EFXo} for instances with
polynomially-bounded weights and small vertex cover, and paint a complete picture
on the parameterized complexity of {\EFXo} parameterized by treewidth, arguably the most well-studied structural parameter.

We refer to \cref{table:algo,table:hardness} for a summary of our algorithms and hardness results respectively.
%\manolis{Perhaps completely avoid the third row? Also, perhaps cite the corresponding theorem for each result.}



\newcommand{\lightrule}{\arrayrulecolor{black!30}\hline\arrayrulecolor{black}}
{
\begin{table}[ht]
\centering
\begin{tabular}{|c|c!{\color{black!30}\vrule}c|}
\hline
& \textbf{Running time} & \textbf{Conditions} \\
\hline
 EF & $\bO(n+m)$ & Binary valuations \\
\lightrule
\multirow{2}{*}{EF} & \multirow{2}{*}{$2^k n^{\bO(1)}$} & At most $k$ heavy edges \\  &  & Simple graph \\
\lightrule

\multirow{2}{*}{EF(X)} & \multirow{2}{*}{$W^{\bO(\tw)}(n+m)^{\bO(1)}$}
 & $W$: \emph{max. shared weight}\\ & &  $\tw$: treewidth\\
\lightrule

\multirow{2}{*}{EF(X)} & \multirow{2}{*}{$(n+m)^{\bO(\tw)}$}
 & Poly weights \\ &  & $\tw$: treewidth \\

\hline
\end{tabular}
\caption{Our algorithms for EF and EFX \textsc{Orientation} in (multi)graphs with $n$ vertices and $m$ edges (\cref{subsec:EFo:linear,subsec:EFo_heavy,subsec:EFo_shared_weight,subsec:EFXo:tw}). All algorithms except the one in the second row also extend to the respective minimum-charity problem.}
\label{table:algo}
\end{table}
}



\begin{table}[ht]
\centering
\begin{tabular}{|c|c!{\color{black!30}\vrule}c|}
\hline
& \textbf{Hardness/} & \multirow{2}{*}{\textbf{Conditions}} \\
& \textbf{Lower Bound} & \\
\hline
 \multirow{2}{*}{EF} & no $(2-\varepsilon)^{k}n^{\bO(1)}$ time algorithm & At most $k$ heavy edges \\ & (under SETH) & Polynomial weights \\
\lightrule
\multirow{2}{*}{EF} & \multirow{2}{*}{NP-complete} & 3-regular graphs \\ & & Weights in \{1,2\}    \\
\lightrule
\multirow{2}{*}{EF} & NP-complete & Vertex cover 2 \textbf{or}\\
& (weakly) & Multigraph with 2 vertices    \\
\lightrule
\multirow{2}{*}{EFX} & NP-complete & Vertex cover 4 \textbf{or}\\
& (weakly) & Multigraph with 4 vertices    \\
\lightrule
\multirow{4}{*}{EF(X)} & W[1]-hard & Parameter: vertex cover number \\
% \lightrule
\arrayrulecolor{black!30}\cline{2-3}\arrayrulecolor{black}
 & XNLP-hard & Parameter: pathwidth   \\
% \lightrule
\arrayrulecolor{black!30}\cline{2-3}\arrayrulecolor{black}
 & XALP-hard & Parameter: treewidth \\ \arrayrulecolor{black!30}\cline{2-3}\arrayrulecolor{black} & & \textbf{All three:} Polynomial weights   \\
\lightrule
\multirow{2}{*}{EF(X)} & no $n^{o(\vc)}$ time algorithm & $\vc$: vertex cover number \\ & (under ETH) & Polynomial weights \\
\hline
\end{tabular}
\caption{Our hardness results for EF and EFX \textsc{Orientation} (\cref{subsec:EF:hardness,subsec:EF->EFX} respectively). All of these results hold even for \textbf{simple symmetric} instances, except the ``multigraph'' results in the third/fourth rows, which are restricted just to symmetric instances.}
\label{table:hardness}
\end{table}
